

%%%%%%%%%%%%%%%%%%%%%%%%%%%%%%%%%%%%%%%%%%%%%%%%%%%%%%%%%%%%%% КОНТРОЛЬ ЦЕПЕЙ ВЫКЛЮЧАТЕЛЯ %%%%%%%%%%%%%%%%%%%%%%%%%%%%%%%%%%%%%%%%%%%%%%%%

\begin{flushleft}
\color{uniblue}{\subsubsection{Контроль цепей выключателя (КЦВ)}}
\color{black}
\nopagebreak
\end{flushleft}

\par
Алгоритм работы функции <<КЦВ>> выполнен в соответствии с рисунком \ref{pdska:img_kzv}.

\vspace{1mm}
\begin{figure}[H]
\centering
\includegraphics[width=0.8\textwidth,height=0.8\textheight,keepaspectratio]{\fbpath/27. ПДС КА/_latex/img/kzv.pdf}
\captionof{figure}{Функциональная схема <<КЦВ>>}\label{pdska:img_kzv}
\end{figure}

%%%%%%%%%%%%%%%%%%%%%%%%%%%%%%%%%%%%%%%%%%%%%%%%%%%%%%%%%% КОНТРОЛЬ ПРУЖИНЫ ВЫКЛЮЧАТЕЛЯ %%%%%%%%%%%%%%%%%%%%%%%%%%%%%%%%%%%%%%%%%%%%%%%%%%%%%

\begin{flushleft}
\color{uniblue}{\subsubsection{Контроль пружины (КПруж)}}
\color{black}
\nopagebreak
\end{flushleft}

\par
Алгоритм работы функции <<КПруж>> выполнен в соответствии с рисунком \ref{pdska:img_kp}.

\vspace{1mm}
\begin{figure}[H]
\centering
\includegraphics[width=0.8\textwidth,height=0.8\textheight,keepaspectratio]{\fbpath/27. ПДС КА/_latex/img/kp.pdf}
\captionof{figure}{Функциональная схема <<КПруж>>}\label{pdska:img_kp}
\end{figure}

%%%%%%%%%%%%%%%%%%%%%%%%%%%%%%%%%%%%%%%%%%%%%%%%%%%%%%%%%% КОНТРОЛЬ ИЗОЛЯЦИИ ВЫКЛЮЧАТЕЛЯ %%%%%%%%%%%%%%%%%%%%%%%%%%%%%%%%%%%%%%%%%%%%%%%%%%

\begin{flushleft}
\color{uniblue}{\subsubsection{Контроль изоляции (КИ)}}
\color{black}
\nopagebreak
\end{flushleft}

\par
Алгоритм работы функции <<КИ>> выполнен в соответствии с рисунком \ref{pdska:img_ki}.

\vspace{1mm}
\begin{figure}[H]
\centering
\includegraphics[width=0.8\textwidth,height=0.8\textheight,keepaspectratio]{\fbpath/27. ПДС КА/_latex/img/ki.pdf}
\captionof{figure}{Функциональная схема <<КИ>>}\label{pdska:img_ki}
\end{figure}

